\documentclass[conference]{IEEEtran}
\IEEEoverridecommandlockouts

\usepackage{cite}
\usepackage{amsmath,amssymb,amsfonts}
\usepackage{graphicx}
\usepackage{textcomp}
\usepackage{booktabs}
\usepackage{multirow}
\usepackage{array}
\usepackage{url}

\begin{document}

\title{Digital Twin Benchmarking of Microclimate Control for
IoT Hydroponic Grow-Boxes: A Two-Period Evaluation of Threshold,
Correlation-Aware, and Q-Learning Strategies}

\author{
\IEEEauthorblockN{Altynay Yertay\IEEEauthorrefmark{1},
Kuanysh Bakirov\IEEEauthorrefmark{2},
Aruzhan Shoman\IEEEauthorrefmark{1}}
\IEEEauthorblockA{\IEEEauthorrefmark{1}Astana IT University, Mangilik El Avenue 55/11, Astana 010000, Kazakhstan}
\IEEEauthorblockA{\IEEEauthorrefmark{2}L.N. Gumilyov Eurasian National University, Satpayev 2, Astana 010008, Kazakhstan}
\IEEEauthorblockA{Corresponding author: k.bakirov@enu.kz, Tel.: +7\,707\,331\,6325}
}

\maketitle

\begin{abstract}
Reinforcement-learning (RL) controllers for agricultural microclimate are
routinely reported to save energy, but such claims are almost always established
on data from the same operating period used for training, so their robustness to
a change of regime is untested. We present a four-layer Digital Twin (DT)
framework for an ESP32-S3-instrumented hydroponic grow-box and use it as a
reproducible harness to benchmark a fixed-threshold baseline, a
correlation-aware orchestrator, and a tabular Q-learning agent with a
comfort-band reward on $\approx$1.3 million telemetry records spanning two
operating periods, over $N=10$ random seeds. In the warmer in-period regime,
where the baseline ventilates on 89.5\% of steps, the agent cuts fan activation
by $49.7\pm18.6\%$ and total actuator energy by $26.0\pm21.9\%$ (Wilcoxon
$p=0.010$). In a held-out future period $\approx$2.5\,$^{\circ}$C cooler, where
the baseline ventilates on only 27.1\% of steps, the frozen policy loses the
advantage entirely: fan activation reaches parity ($p=0.49$) while total energy
\emph{rises} by $95.8\pm65.9\%$ ($p=0.002$), driven by 17.4\% residual heater
duty. A sensitivity analysis over the comfort band, the energy weight and the
telemetry-grounding ratio reproduces this reversal in all eleven configurations
tested, including pure telemetry replay with the simulation layer disabled, so it
is a property of the regime rather than of the reward tuning or the Twin. From
the energy balance we derive a break-even condition: out-of-period the agent
would have to remove 82.7 percentage points of fan duty to pay for its own
heating---more than the baseline uses in total. The $16.7$\,KB Q-table runs on
the ESP32-S3 unaided. The contribution is a reproducible DT benchmark together
with a quantitative criterion for when RL is worth deploying.
\end{abstract}

\begin{IEEEkeywords}
Benchmarking; Digital Twin; Edge Computing; ESP32; Generalisation;
Hydroponic Microclimate Control; IoT; Precision Agriculture; Q-Learning;
Reinforcement Learning; Vertical Farming
\end{IEEEkeywords}

\section{Introduction}
Controlled Environment Agriculture and urban vertical farming are emerging as
sustainable responses to global food-security pressures, with hydroponic systems
reducing water consumption by up to 95\% compared with soil-based cultivation
\cite{kozai2022,oh2023}. In semi-arid regions such as Kazakhstan, small-scale
grow-boxes are an attractive food-resilience strategy, and the quality of
automated microclimate control directly determines yield, operational cost, and
energy footprint \cite{wolfert2017,sneineh2023}.

Contemporary IoT-enabled grow-boxes rely primarily on static threshold rules,
which treat each sensor channel independently and can waste energy by triggering
actuators unnecessarily \cite{ferentinos2017,tzounis2017}. Reinforcement learning
(RL) is widely proposed as a remedy, with reported energy savings of 7--57\%
against rule-based baselines \cite{adesanya2024,ajagekar2022}. A recurring
methodological weakness, however, is that these savings are typically measured on
held-out data from the \emph{same} period and installation used for training;
whether a learned policy still helps once the operating regime shifts---when
ambient temperature drifts and the baseline is no longer over-reactive---is
rarely tested. For resource-constrained deployments, where retraining is costly
and policies are often frozen on-device, this question is decisive.

Digital Twin (DT) frameworks allow it to be studied offline, without disrupting
live cultivation. Verdouw et al.\ \cite{verdouw2021} provide a foundational
IoT-A-based DT framework validated across five agricultural use cases; Luo and
Ball \cite{luo2025} apply Q-learning within a DT harness, but for weekly
production scheduling rather than second-by-second microclimate regulation on
real telemetry. We build on this line of work with a benchmarking focus and make
three contributions:
\begin{enumerate}
\item A four-layer DT framework synchronised with $\approx$1.3 million real
ESP32-S3 telemetry records across \emph{two} operating periods, providing a
reproducible offline harness for benchmarking control algorithms.
\item A two-period, multi-seed evaluation protocol---an in-period chronological
split and a genuinely \emph{future}, continuous hold-out in a cooler
regime---applied to three controllers under identical conditions, with a
sensitivity analysis over the reward and simulation design parameters.
\item A quantitative characterisation of when RL helps: in-period gains
($-49.7\pm18.6\%$ fan, $-26.0\pm21.9\%$ energy) do not transfer (fan parity and
$+95.8\pm65.9\%$ energy out-of-period), together with a closed-form break-even
condition that predicts the sign of the outcome from baseline telemetry alone,
and an edge-deployable $16.7$\,KB Q-table.
\end{enumerate}

\section{Related Work}
\textbf{Digital Twins in agriculture.}
Verdouw et al.\ \cite{verdouw2021} define DT types and an IoT-A implementation
model validated across five IoF2020 use cases, forming the conceptual basis for
our four-layer architecture. Ariesen-Verschuur et al.\ \cite{ariesen2022} review
123 studies on DT in greenhouse horticulture, identifying real-time sensor
synchronisation and physics-based simulation as key enablers and noting that
integrating adaptive control within DT frameworks remains open. Bua et al.\
\cite{bua2025} apply an RL-driven DT to greenhouse robotics under communication
latency; Luo and Ball \cite{luo2025} apply a prescriptive DT to vertical-farm
scheduling on historical demand rather than live telemetry. We use the DT
specifically as a reproducible \emph{benchmarking harness} and add an explicit
temporal-generalisation test.

\textbf{RL for climate control.}
Adesanya et al.\ \cite{adesanya2024} apply deep Q-learning to PID tuning in a
TRNSYS-Python cosimulation, reporting 7.17--12.99\% energy reductions; Ajagekar
and You \cite{ajagekar2022} report $\approx$57\% reduction with a deep Q-network
in semi-closed greenhouses. Chen et al.\ \cite{chen2025} survey greenhouse
control and note that RL is gaining traction for adapting without an explicit
plant model. Two gaps motivate our study. First, computational footprint: deep
Q-networks require compute unavailable on ESP32-class microcontrollers, whereas
our tabular agent targets on-device execution. Second, evaluation discipline:
reported savings are seldom tested across a genuine regime shift, which is what
our held-out period probes. Table~\ref{tab:related} positions this study against
the closest prior work: to our knowledge none combines a DT harness, an RL
controller, microcontroller-class deployment, and a regime-shift evaluation.

\textbf{IoT-based hydroponic monitoring.}
Sneineh and Shabaneh \cite{sneineh2023} show that ESP32 microcontrollers with pH,
TDS, temperature, and water-level sensors suffice for grow-box monitoring;
Rofiansyah et al.\ \cite{rofiansyah2025} add CNN-based plant-health assessment.
Wolfert et al.\ \cite{wolfert2017} identify real-time telemetry integration as
the foundational challenge, and Ferentinos et al.\ \cite{ferentinos2017} confirm
that short sampling intervals capture actuation-relevant dynamics. Our platform
extends monitoring toward closed-loop adaptive control evaluated through the DT
harness.

\begin{table}[t]
\caption{Positioning against related DT/RL work for controlled-environment agriculture}
\label{tab:related}
\centering
\footnotesize
\setlength{\tabcolsep}{3pt}
\begin{tabular}{lcccc}
\toprule
Work & DT & RL & MCU-class & Regime-shift \\
     &    &    & deployment & evaluation \\
\midrule
Verdouw et al.\ \cite{verdouw2021}     & \checkmark & --         & --         & -- \\
Ariesen-V.\ et al.\ \cite{ariesen2022} & \checkmark & --         & --         & -- \\
Luo and Ball \cite{luo2025}            & \checkmark & \checkmark & --         & -- \\
Bua et al.\ \cite{bua2025}             & \checkmark & \checkmark & --         & -- \\
Adesanya et al.\ \cite{adesanya2024}   & --         & \checkmark & --         & -- \\
Ajagekar and You \cite{ajagekar2022}   & --         & \checkmark & --         & -- \\
Sneineh and Sh.\ \cite{sneineh2023}    & --         & --         & \checkmark & -- \\
\midrule
This work & \checkmark & \checkmark & \checkmark & \checkmark \\
\bottomrule
\end{tabular}
\end{table}

\section{System Architecture}

\subsection{Physical Grow-Box}
Experiments were carried out on a closed, insulated hydroponic cabinet of roughly
0.8\,m\textsuperscript{3} in the AgroTech laboratory at Astana IT University. It
holds four growing tiers of microgreens and leafy vegetables and is actuated by
LED grow lighting, a speed-controlled ventilation fan, a resistive heater, and a
circulating nutrient pump. An ESP32-S3 reads fourteen channels in four groups
\cite{bakirov2025}: an SHT40 and an SCD41 for air temperature, humidity and
CO\textsubscript{2} (agro-climate); pH, EC, TDS, turbidity and water-temperature
probes (hydro-chemistry); a BH1750 illuminance sensor; and five capacitive
substrate-moisture probes (root zone). Samples are taken every $\approx$5.6\,s
and pushed as JSON over HTTP to an on-site Python/Flask gateway
\cite{akhmetzhanov2024,yedilkhan2025}, which writes one sixteen-column CSV file
per day. Across the trial only the three agro-climate variables and illuminance
returned dependable values, and these define the control state; the
hydro-chemistry and root-zone channels were faulty (Sec.~\ref{sec:limitations}).

\subsection{Dataset: Two Operating Periods}
Telemetry was logged from a single ESP32-S3 node. Collection was intermittent
during platform development: the full sixteen-column schema stabilised only
toward the end of May 2025, and logging was interrupted for extended stretches,
most notably a four-month outage from October 2025 to January 2026. The two
periods were archived separately and are non-overlapping by construction; the
loader additionally verifies that no timestamp occurs in both archives, so
nothing from the hold-out can leak into training.

\emph{Period 1 (in-period), 29 May 2025--28 April 2026:} 137 daily files,
808{,}078 raw records, 650{,}126 after quality filtering ($T \notin
[15,40]\,^{\circ}$C; $\mathit{RH} \notin [10,95]\%$; $\mathrm{CO_2} \notin
[300,2000]$\,ppm). The data are sparse---present on only 100 of 334 calendar days
($\approx$12.6\% duty), with $\approx$85\% of records in February--April 2026,
immediately before the hold-out. They were split chronologically 80/20 into
520{,}100 training (to 27 March 2026) and 130{,}026 in-period evaluation samples.

\emph{Period 2 (held-out), 29 April--15 June 2026:} 48 daily files, 726{,}781 raw
records, 652{,}909 after filtering---a genuinely future, \emph{continuous} block
(44 of 48 days, $\approx$89\% duty) used exclusively for evaluation. Since
Period 1 is concentrated in early 2026 and Period 2 immediately follows it, this
is a \emph{near-term forward-transfer} test on one installation rather than a
cross-seasonal one.

Table~\ref{tab:env} contrasts the full periods with the two 60{,}000-step
segments actually evaluated; the latter, not the former, generate the results.
The operative difference between the segments is \emph{thermal}, not
hygrometric: mean humidity is nearly identical (53.5\% versus 55.0\%) while
temperature drops by $\approx$2.5\,$^{\circ}$C. Since the threshold baseline
ventilates on $T>25.5\,^{\circ}$C, the decisive statistic is the fraction of
steps above that threshold: 89.5\% in-period against 27.1\% out-of-period. This
single number, visible in the telemetry before any controller is run, accounts
for the entire result of Sec.~\ref{sec:generalization}. Mean
CO\textsubscript{2} stays well below 800\,ppm throughout, confirming that fan
demand is thermal rather than CO\textsubscript{2}-driven.

\begin{table}[t]
\caption{Environment statistics (mean $\pm$ SD): full quality-filtered periods
versus the two evaluated 60{,}000-step windows}
\label{tab:env}
\centering
\footnotesize
\setlength{\tabcolsep}{3pt}
\begin{tabular}{lcccc}
\toprule
 & \multicolumn{2}{c}{Full period} & \multicolumn{2}{c}{Evaluation window} \\
\cmidrule(lr){2-3}\cmidrule(lr){4-5}
Parameter & P1 & P2 & P1 & P2 \\
\midrule
Temperature ($^{\circ}$C)  & $26.4\pm1.9$ & $25.1\pm1.7$ & $26.8\pm1.1$ & $24.3\pm1.5$ \\
Humidity (\%)              & $39.5\pm10.5$ & $58.0\pm8.3$ & $53.5\pm5.1$ & $55.0\pm10.4$ \\
CO\textsubscript{2} (ppm)  & $569\pm183$ & $525\pm130$ & --- & --- \\
Steps $>25.5\,^{\circ}$C (\%) & --- & --- & 89.5 & 27.1 \\
\bottomrule
\end{tabular}
\end{table}

\subsection{Digital Twin Architecture}
The DT follows the IoT-A reference structure \cite{verdouw2021} with four layers:
\textit{Perception} (ESP32-S3 sensors); \textit{Edge} (Python/Flask gateway
serialising JSON telemetry to CSV); \textit{Twin} (offline Python simulation
engine reading the CSV archive and applying a physics-based actuator-effect
model); and \textit{Control} (a pluggable interface accepting
$(T,\mathit{RH},\mathrm{CO_2},\text{light},\text{history})$ and returning binary
actuator commands). The physics model estimates per-step effects: the fan reduces
temperature by 0.10\,$^{\circ}$C, humidity by 0.6 points and
CO\textsubscript{2} by 18\,ppm; the heater raises temperature by
0.15\,$^{\circ}$C and lowers humidity by 0.3 points; the pump raises humidity by
0.9 points. Ambient drift is a mean-reversion term whose coefficient and target
are identified per period by least squares on consecutive telemetry samples, the
residual standard deviation supplying the Gaussian measurement noise; the
resulting one-step prediction error is 0.005\,$^{\circ}$C in-period and
0.003\,$^{\circ}$C out-of-period. During evaluation the Twin blends real
telemetry and simulated state at a 70/30 ratio every step (60/40 every 50 steps
during training), grounding the rollout in measured conditions while preserving
causal sensitivity to distinguish policies \cite{luo2025}. This harness is the
core contribution: any controller can be evaluated on identical real-world
sequences, reproducibly and without affecting live plants.

\section{Control Algorithms}

\subsection{Threshold Baseline}
Static hysteresis rules act independently per channel: the fan activates when
temperature exceeds 25.5\,$^{\circ}$C or CO\textsubscript{2} exceeds 800\,ppm;
the heater below 20.0\,$^{\circ}$C; the pump runs continuously. This is the
predominant strategy for low-cost hydroponic systems
\cite{sneineh2023,rofiansyah2025}.

\subsection{Correlation-Aware Orchestrator}
The orchestrator adds a proactive layer: when relative humidity rises by more
than 2.0 points over 8 steps and temperature simultaneously exceeds
23.5\,$^{\circ}$C, the fan is pre-activated before the 25.5\,$^{\circ}$C
threshold is reached, exploiting the observed lag between
evapotranspiration-driven humidity rise and subsequent temperature rise
\cite{bakirov2025,ferentinos2017}. All other rules match the baseline.

\subsection{Q-Learning Agent}
The agent discretises the three primary channels into $7\times7\times5$ bins plus
a binary light indicator (lux $>50$), yielding 490 states, and chooses among 8
binary combinations of fan, heater and pump. Penalising every deviation from a
single 24\,$^{\circ}$C setpoint would, in an enclosure whose mean temperature
sits above the setpoint, drive continuous over-cooling followed by compensatory
heating---a prediction we verify empirically in Sec.~\ref{sec:sensitivity}. We
therefore use a \emph{comfort-band} temperature term, zero inside
$[T_{\text{low}},T_{\text{high}}]=[22,27]\,^{\circ}$C and quadratic outside,
reflecting the achievable operating range under sustained LED load. With
$d_T=\max(0,\,T_{\text{low}}-T,\,T-T_{\text{high}})$,
$e_{\mathit{RH}}=|\mathit{RH}-55|$ and $[\cdot]^{+}=\max(0,\cdot)$:
\begin{equation}
r = -\!\left(\frac{d_T^{2}}{16}
+ 0.5\,\frac{e_{\mathit{RH}}^{2}}{400}
+ 0.3\,\frac{[\mathrm{CO_2}-800]^{+}}{200}
+ 0.04\,E_{\text{energy}}\right)
\end{equation}
where
$E_{\text{energy}}=1.0\cdot\mathbf{1}_{\text{fan}}+1.5\cdot\mathbf{1}_{\text{heater}}+0.3\cdot\mathbf{1}_{\text{pump}}$
encodes relative power draw, and a safety penalty $r=-25$ applies outside
$[16,30]\,^{\circ}$C \cite{ajagekar2022,chen2025}. Q-values follow the standard
Bellman update with $\alpha=0.15$, $\gamma=0.95$ and an $\varepsilon$-greedy
schedule decaying geometrically (factor 0.9997 per step) to a floor of 0.05
within the first of 8 training episodes over the 520{,}100-step training set.
Training updates $267.2\pm4.4$ of the 490 states; states never updated carry no
learned policy and the controller falls back to the threshold rules there, so the
deployed policy is never undefined.

\section{Experimental Setup}
All controllers are evaluated on the same 60{,}000-step chronological segment
within each period ($\approx$94 hours at the native cadence), receiving identical
real readings as the 70\% grounding signal so that any outcome difference is
attributable to the control policy alone. We report temperature and humidity RMSE
(relative to 24\,$^{\circ}$C and 55\%), fan, heater and pump activation rates,
and the microclimate excursion rate (fraction of steps outside
$[18,28]\,^{\circ}$C). Total actuator energy assumes 18\,W (fan), 120\,W (heater)
and 10\,W (pump) over a 730-hour month.

Stochasticity enters at three points: the $\varepsilon$-greedy exploration and
Q-table initialisation during training, and the Gaussian noise term of the Twin
during evaluation. We therefore repeat the full train--evaluate pipeline for
$N=10$ seeds and report mean $\pm$ standard deviation throughout. Within a seed
all three controllers are driven by the \emph{same} pre-drawn noise stream, so
comparisons are paired; significance of the Q-learning--threshold difference is
assessed with a two-sided Wilcoxon signed-rank test over the $N$ paired runs.

\section{Results and Discussion}

\subsection{In-Period Comparison}
Table~\ref{tab:multiseed} reports both periods. In Period 1 the Q-learning agent
reduces fan activation from 89.7\% to $45.1\pm16.7\%$ (a $49.7\pm18.6\%$ relative
reduction, $p=0.002$) and pump activation from 100\% to $68.1\pm8.2\%$, while
holding temperature RMSE within 0.07\,$^{\circ}$C of the baseline and lowering
humidity RMSE from 5.25 to $4.97\pm0.13$. The microclimate excursion rate rises
from 7.00\% to $8.33\pm1.36\%$---the price paid for the energy reduction, still
far from the safety band.

\begin{table*}[t]
\caption{Multi-seed comparison (mean $\pm$ SD over $N=10$ seeds,
60{,}000-step evaluation window in each period)}
\label{tab:multiseed}
\centering
\begin{tabular}{llcccccc}
\toprule
Period & Algorithm & RMSE-T ($^{\circ}$C) & RMSE-RH (\%) & Fan\% & Heat\% & Pump\% & Energy (kWh/mo) \\
\midrule
\multirow{3}{*}{P1 (in-period)}
 & Threshold    & $2.89\pm0.00$ & $5.25\pm0.00$ & $89.7\pm0.0$ & $0.0\pm0.0$ & $100.0\pm0.0$ & $19.09\pm0.00$ \\
 & Correlation  & $2.89\pm0.00$ & $5.25\pm0.00$ & $89.7\pm0.0$ & $0.0\pm0.0$ & $100.0\pm0.0$ & $19.09\pm0.00$ \\
 & Q-Learning   & $2.96\pm0.03$ & $4.97\pm0.13$ & $45.1\pm16.7$ & $3.7\pm3.9$ & $68.1\pm8.2$ & $14.12\pm4.18$ \\
\midrule
\multirow{3}{*}{P2 (held-out)}
 & Threshold    & $1.49\pm0.00$ & $10.42\pm0.00$ & $32.3\pm0.0$ & $0.0\pm0.0$ & $100.0\pm0.0$ & $11.55\pm0.00$ \\
 & Correlation  & $1.49\pm0.00$ & $10.42\pm0.00$ & $32.4\pm0.0$ & $0.0\pm0.0$ & $100.0\pm0.0$ & $11.55\pm0.00$ \\
 & Q-Learning   & $1.53\pm0.02$ & $9.95\pm0.12$ & $33.8\pm8.4$ & $17.4\pm7.5$ & $39.7\pm5.2$ & $22.61\pm7.60$ \\
\bottomrule
\end{tabular}
\\[2pt]
\footnotesize The threshold and correlation controllers are deterministic given
the telemetry; their residual seed variance comes only from the Twin noise term
and is below the reported precision on every metric.
\end{table*}

\subsection{Temporal Generalisation (Held-Out Period)}
\label{sec:generalization}
The central result concerns transfer to the held-out period, which
Fig.~\ref{fig:gen} contrasts with the in-period outcome. On Period 2 the
advantage \emph{disappears}: fan activation ($33.8\pm8.4\%$) is statistically
indistinguishable from the threshold baseline's (32.3\%; $p=0.49$), so the agent
provides no cooling-energy benefit. The mechanism is structural rather than
stochastic and is visible in the telemetry before any controller is run: the
baseline ventilates on 89.5\% of steps in the in-period window but on only 27.1\%
in the held-out window (Table~\ref{tab:env}), because the latter is
$\approx$2.5\,$^{\circ}$C cooler. There is simply no over-ventilation left to
remove. Microclimate accuracy is high for all controllers there (RMSE-T
$\approx$1.5\,$^{\circ}$C, zero excursions) because the cooler regime sits closer
to the setpoint. Humidity RMSE is larger ($\approx$10.4\% versus $\approx$5.3\%)
for every controller alike; this reflects the greater variability of the held-out
window (RH standard deviation 10.4\% versus 5.1\%) rather than any control
failure, since its mean humidity sits exactly at the reference.

\begin{figure}[t]
\centering
\includegraphics[width=\columnwidth]{generalization.png}
\caption{Total actuator energy, mean $\pm$ SD over 10 seeds. In Period 1 the
agent saves $26.0\%$; in the held-out Period 2 it costs $95.8\%$ \emph{more}
than the baseline, driven by compensatory heating.}
\label{fig:gen}
\end{figure}

\subsection{Energy Accounting Is Regime-Dependent}
In Period 1 total actuator energy falls from 19.09 to $14.12\pm4.18$\,kWh/month
($-26.0\pm21.9\%$, $p=0.010$): the fan and pump reductions outweigh the agent's
3.7\% heater duty. The spread deserves emphasis---eight of ten seeds save energy
and two do not---so the direction is reliable while the magnitude is not. In
Period 2 total energy \emph{rises} from 11.55 to $22.61\pm7.60$\,kWh/month
($+95.8\pm65.9\%$, $p=0.002$). The cooler regime pushes the state into bins whose
learned action prescribes heating: heater duty climbs from 3.7\% to
$17.4\pm7.5\%$, and at 120\,W---nearly seven times the fan's draw---this
dominates the balance even though pump duty falls to 39.7\%.

Because all three actuators are binary, the comparison reduces to a closed form.
With duty fractions $f,h,p$ the mean power is $P=P_f f+P_h h+P_p p$; the baseline
has $h=0$ and $p=1$, so the agent saves energy if and only if
\begin{equation}
\Delta f \;>\; \Delta f^{*} \;=\;
\frac{P_h\,h_{\text{QL}} + P_p\,(p_{\text{QL}}-1)}{P_f},
\label{eq:breakeven}
\end{equation}
where $\Delta f = f_{\text{th}} - f_{\text{QL}}$ is the fan duty the agent
removes. In-period, $\Delta f^{*}=6.8$ percentage points against an observed
$44.5$, hence the saving. Out-of-period, $\Delta f^{*}=82.7$ against an observed
$-1.5$: the requirement exceeds the baseline's entire fan duty (32.3\%), so the
deficit cannot be closed by \emph{any} reduction in ventilation whatsoever.
Equation~\eqref{eq:breakeven} needs only a few days of baseline telemetry
together with the agent's heater duty, and can therefore be evaluated before
deployment. The practical lesson is that headline ``RL energy savings'' can be an
artifact of an over-reactive baseline under high thermal load, and that
total-system accounting must include compensatory actuation, not fan energy
alone.

\subsection{Sensitivity to Reward and Simulation Design}
\label{sec:sensitivity}
Table~\ref{tab:sensitivity} varies each design parameter in turn, holding the
others at their default. The comfort band is decisive in-period: replacing it
with a single 24\,$^{\circ}$C setpoint reproduces exactly the pathology it was
introduced to prevent, turning the $-27.0\%$ in-period saving into a $+26.2\%$
penalty, and narrowing the band to $[23,26]$ almost erases the saving ($+4.8\%$).
Widening it to $[21,28]$ improves the in-period result to $-38.0\%$, and a
heavier energy weight helps monotonically ($-35.2\%$ at $w=0.08$).

None of this rescues the held-out period. All eleven configurations show an
energy increase there, from $+46.3\%$ to $+160.5\%$; even the best in-period
setting, $[21,28]$, still costs $+80.4\%$ out-of-period, so the reversal is not
an artifact of how the reward was tuned. Varying the grounding ratio changes the
outcome by less than five percentage points across the whole range, and the
$100/0$ condition---pure telemetry replay with the simulated dynamics switched
off entirely---gives $-24.7\%$ and $+89.4\%$, within a point of the default. The
central finding is therefore not an artifact of the simulation layer either.

\begin{table}[t]
\caption{Sensitivity of the Q-learning agent (mean over $N=5$ seeds;
$\Delta E$ relative to the threshold baseline of the same period)}
\label{tab:sensitivity}
\centering
\footnotesize
\setlength{\tabcolsep}{3pt}
\begin{tabular}{llcccc}
\toprule
\multirow{2}{*}{Parameter} & \multirow{2}{*}{Value} &
\multicolumn{2}{c}{Period 1} & \multicolumn{2}{c}{Period 2} \\
\cmidrule(lr){3-4}\cmidrule(lr){5-6}
 & & Fan\% & $\Delta E$\% & Fan\% & $\Delta E$\% \\
\midrule
\multirow{4}{*}{Comfort band}
 & single 24\,$^{\circ}$C & 86.0 & $+26.2$ & 49.9 & $+160.5$ \\
 & $[21,28]$ & 44.6 & $-38.0$ & 28.4 & $+80.4$ \\
 & $[22,27]^{\dagger}$ & 51.4 & $-27.0$ & 33.1 & $+84.3$ \\
 & $[23,26]$ & 73.0 & $+4.8$ & 48.0 & $+124.6$ \\
\midrule
\multirow{3}{*}{Energy weight}
 & 0.02 & 55.6 & $-16.3$ & 30.6 & $+99.7$ \\
 & $0.04^{\dagger}$ & 51.4 & $-27.0$ & 33.1 & $+84.3$ \\
 & 0.08 & 47.5 & $-35.2$ & 26.5 & $+46.3$ \\
\midrule
\multirow{4}{*}{Grounding ratio}
 & 100/0 & 53.8 & $-24.7$ & 36.3 & $+89.4$ \\
 & 80/20 & 51.8 & $-26.4$ & 34.2 & $+87.5$ \\
 & $70/30^{\dagger}$ & 51.4 & $-27.0$ & 33.1 & $+84.3$ \\
 & 60/40 & 51.0 & $-27.5$ & 32.0 & $+76.7$ \\
\bottomrule
\end{tabular}
\\[2pt]
\footnotesize $^{\dagger}$ default configuration. Its Fan\% differs from
Table~\ref{tab:multiseed} because the sensitivity sweep uses the first five
seeds rather than all ten; averaged over those same five seeds the two agree.
\end{table}

\subsection{Orchestrator and Edge Deployment}
The orchestrator's proactive rule changed the commanded action on 0.023\% of
steps in Period 1 and 0.049\% in Period 2. Because all controllers share an
identical noise stream within a seed, the paired design detects even this: the
fan-duty difference is nominally significant ($p=0.002$ in both periods), yet it
amounts to $+0.023$ and $+0.047$ percentage points and no other metric moves at
the reported precision---a case where statistical significance and practical
relevance come apart. The humidity$\rightarrow$temperature coupling is physically
real but far too rare at actuation scale to matter here; the held-out segment
(55.0\% RH) is also barely more humid than the in-period one (53.5\%), so it is a
weak test of the mechanism.

Independent of the energy outcome, the tabular agent is edge-viable: training
updates $267.2\pm4.4$ of the 490 states, giving a Q-table of $16.7\pm0.3$\,KB,
three orders of magnitude below the 512\,KB RAM and 4\,MB flash of the ESP32-S3.
The remaining 223 states carry no learned policy; the firmware falls back to the
threshold rules there, so the controller degrades gracefully to the baseline
rather than to an arbitrary action in unseen conditions. The table can be
serialised as a C array and compiled into firmware for near-zero-cost inference,
unlike deep RL approaches \cite{ajagekar2022,bua2025} that require dedicated
compute; even neural models now reach microcontrollers through TinyML toolchains
\cite{ray2022}, but a tabular policy sidesteps that machinery altogether.

\subsection{Limitations}
\label{sec:limitations}
(i)~Direction and magnitude have different reliability: the held-out energy
increase appears in ten of ten seeds and in all eleven sensitivity
configurations, but the in-period energy saving holds in only eight of ten
($-26.0\pm21.9\%$), so single-seed figures of the kind commonly reported in this
literature can mislead by a wide margin in either direction. (ii)~The sensitivity
analysis varies parameters individually; interactions between them, and the
state-discretisation resolution, were not swept. (iii)~The Twin's actuator model
is linear and omits nonlinear heat transfer and latent-heat coupling. This is a
deliberate trade-off: the drift coefficients are identified directly from logged
telemetry (one-step prediction error 0.005\,$^{\circ}$C in-period,
0.003\,$^{\circ}$C out-of-period) and the 70/30 grounding re-anchors the rollout
to measured conditions at every step, bounding accumulated drift; the close
agreement between the $100/0$ and $70/30$ conditions
(Table~\ref{tab:sensitivity}) indicates that this simplification does not drive
the conclusions. The Twin is adequate as a \emph{comparative} harness---all
policies face identical dynamics---but is not a predictive simulator, and
absolute energy figures should be read as estimates. (iv)~The data are
intermittent and concentrated in early 2026, so the held-out comparison is a
near-term forward transfer on a single installation rather than a cross-seasonal
one. (v)~Ten of the fourteen channels (EC, TDS, pH, turbidity, water temperature
and five capacitive substrate probes) were excluded from the control state
because of hardware faults: the pH probe returned values outside the physical
0--14 range, indicating a calibration or reference-electrode failure, and the
capacitive probes returned constant readings consistent with a disconnected ADC
channel. The control scope is therefore restricted to agro-climate regulation;
extending it requires probe replacement and recalibration, and is orthogonal to
the temporal-generalisation question studied here. (vi)~The study covers a single
grow-box of fixed geometry; generalisation to other enclosures and crops requires
further work \cite{chen2025}.

\section{Conclusion}
We presented a four-layer Digital Twin framework serving as a reproducible,
non-disruptive harness for benchmarking microclimate control in an IoT-enabled
hydroponic grow-box, and used it to evaluate three controllers across two
operating periods over ten random seeds. Tabular Q-learning is genuinely useful
\emph{when the baseline is over-reactive}, but it is \emph{not} a drop-in frozen
energy optimiser: the in-period gains reversed into a $95.8\pm65.9\%$ energy
increase on a cooler held-out period, and the sensitivity analysis reproduced
that reversal in all eleven configurations tested, showing it to be a property of
the operating regime rather than of the reward design or the simulation layer.
This is consistent with the broader RL literature, in which policies overfit to
their training regime \cite{kirk2023}, and with reviews of RL for building
control that stress the hurdles of real-world deployment \cite{wanghong2020}.
Equation~\eqref{eq:breakeven} turns this into a deployment test evaluable from
baseline telemetry before any training. Future work will add periodic on-device
re-adaptation, evaluate on the most humid portion of Period 2, extend the harness
to multi-zone grow-boxes, and couple the control layer with an LLM-based advisory
interface \cite{bakirov2025}.

\section*{Data and Code Availability}
The DT harness, the three controllers, the multi-seed and sensitivity scripts,
and the quality-filtered telemetry for both periods are available at
\url{https://github.com/REPLACE-ME}.

\section*{Acknowledgment}
This research was funded by the Committee of Science of the Ministry of Science
and Higher Education of the Republic of Kazakhstan (Grant No.\ BR24992852,
``Intelligent models and methods of the Smart City digital ecosystem for
sustainable development and the improvement of the citizens' quality of life'').

\begin{thebibliography}{20}
\bibitem{kozai2022} T. Kozai, G. Niu, and M. Takagaki, Eds., \emph{Plant Factory:
An Indoor Vertical Farming System for Efficient Quality Food Production}, 2nd ed.
London, U.K.: Academic Press, 2020.
\bibitem{oh2023} S. Oh and C. Lu, ``Vertical farming -- smart urban agriculture for
enhancing resilience and sustainability in food security,'' \emph{J. Hortic. Sci.
Biotechnol.}, vol. 98, pp. 133--140, 2023.
\bibitem{wolfert2017} S. Wolfert, L. Ge, C. Verdouw, and M. J. Bogaardt, ``Big data
in smart farming -- a review,'' \emph{Agric. Syst.}, vol. 153, pp. 69--80, 2017.
\bibitem{sneineh2023} A. Abu Sneineh and A. A. Shabaneh, ``Design of a smart
hydroponics monitoring system using an ESP32 microcontroller and the Internet of
Things,'' \emph{MethodsX}, vol. 11, Art. no. 102401, 2023.
\bibitem{ferentinos2017} K. P. Ferentinos, N. Katsoulas, A. Tzounis, T. Bartzanas,
and C. Kittas, ``Wireless sensor networks for greenhouse climate and plant
condition assessment,'' \emph{Biosyst. Eng.}, vol. 153, pp. 70--81, 2017.
\bibitem{tzounis2017} A. Tzounis, N. Katsoulas, T. Bartzanas, and C. Kittas,
``Internet of Things in agriculture, recent advances and future challenges,''
\emph{Biosyst. Eng.}, vol. 164, pp. 31--48, 2017.
\bibitem{adesanya2024} M. Adesanya et al., ``Deep reinforcement learning for PID
parameter tuning in greenhouse HVAC system energy optimization,'' \emph{Expert
Syst. Appl.}, vol. 252, Art. no. 124126, 2024.
\bibitem{ajagekar2022} A. Ajagekar and F. You, ``Deep reinforcement learning based
automatic control in semi-closed greenhouse systems,'' \emph{IFAC-PapersOnLine},
vol. 55, no. 7, pp. 406--411, 2022.
\bibitem{verdouw2021} C. Verdouw, B. Tekinerdogan, A. Beulens, and S. Wolfert,
``Digital twins in smart farming,'' \emph{Agric. Syst.}, vol. 189, Art. no. 103046,
2021.
\bibitem{luo2025} Y. Luo and P. Ball, ``Adaptive production strategy in vertical
farm digital twins with Q-learning algorithms,'' \emph{Sci. Rep.}, vol. 15,
Art. no. 15129, 2025.
\bibitem{ariesen2022} N. Ariesen-Verschuur, C. Verdouw, and B. Tekinerdogan,
``Digital twins in greenhouse horticulture: a review,'' \emph{Comput. Electron.
Agric.}, vol. 199, Art. no. 107183, 2022.
\bibitem{bua2025} C. Bua, L. Borgianni, D. Adami, and S. Giordano, ``Reinforcement
learning-driven digital twin for zero-delay communication in smart greenhouse
robotics,'' \emph{Agriculture}, vol. 15, no. 12, Art. no. 1290, 2025.
\bibitem{chen2025} S. Chen, A. Liu, F. Tang, P. Hou, Y. Lu, and P. Yuan, ``A review
of environmental control strategies and models for modern agricultural
greenhouses,'' \emph{Sensors}, vol. 25, no. 5, Art. no. 1388, 2025.
\bibitem{rofiansyah2025} W. Rofiansyah, F. R. Zalianty, F. A. La Ito, I. Wijayanto,
H. H. Ryanu, and I. D. Irawati, ``IoT-based control and monitoring system for
hydroponic plant growth using image processing and mobile applications,''
\emph{PeerJ Comput. Sci.}, vol. 11, Art. no. e2763, 2025.
\bibitem{bakirov2025} K. Bakirov, A. Kenzhebai, J. Tussupov, I. Shayea, A. Shoman,
and D. Yedilkhan, ``Integrating AI-based monitoring system for microgreen growth
in vertical farming,'' in \emph{Proc. IEEE 5th Int. Conf. Smart Information
Systems and Technologies (SIST)}, 2025, pp. 1--6.
\bibitem{akhmetzhanov2024} B. Akhmetzhanov et al., ``Multi-layer integration of
heterogeneous wireless sensor networks for smart home optimization,'' \emph{Procedia
Comput. Sci.}, vol. 231, pp. 666--671, 2024.
\bibitem{yedilkhan2025} D. Yedilkhan, A. Turginbekov, B. Akhmetzhanov, Z. Talgatuly,
and D. Marlambekov, ``Comparative analysis of IoT wireless technologies
communications through experimental simulations,'' \emph{Procedia Comput. Sci.},
vol. 257, pp. 382--387, 2025.
\bibitem{ray2022} P. P. Ray, ``A review on TinyML: state-of-the-art and prospects,''
\emph{J. King Saud Univ. Comput. Inf. Sci.}, vol. 34, no. 4, pp. 1595--1623, 2022.
\bibitem{kirk2023} R. Kirk, A. Zhang, E. Grefenstette, and T. Rockt\"aschel,
``A survey of zero-shot generalisation in deep reinforcement learning,''
\emph{J. Artif. Intell. Res.}, vol. 76, pp. 201--264, 2023.
\bibitem{wanghong2020} Z. Wang and T. Hong, ``Reinforcement learning for building
controls: the opportunities and challenges,'' \emph{Appl. Energy}, vol. 269,
Art. no. 115036, 2020.
\end{thebibliography}

\end{document}
